\documentclass[11pt, letterpaper]{article}
\input{/home/hyperion/.config/LatexHub/preambles/base.tex}
\input{/home/hyperion/.config/LatexHub/preambles/diagrams.tex}
\input{/home/hyperion/.config/LatexHub/preambles/tikz.tex}
\input{/home/hyperion/.config/LatexHub/preambles/macros-cat-theory.tex}
\input{/home/hyperion/.config/LatexHub/preambles/macros-algebra.tex}
\input{/home/hyperion/.config/LatexHub/preambles/basic-theorems-section.tex}
\input{/home/hyperion/.config/LatexHub/preambles/refs-basic.tex}
\theoremstyle{plain}
\newtheorem{problem}{Problem}
\usepackage{titlepageBU}
\title{Homework 4}
\begin{document}
\makeproblem

\begin{problem}
	Given an example of an open subset $U\subseteq \mathbb{C} $ and two holomorphic functions $f,g : U \to \mathbb{C} $ such that $f'=g'$ but $g$ does not equal $f$ plus a constant.
\end{problem}
\begin{proof}
	Let $U$ be the union of two disjoint open connected subsets $U_1,U_2 \subseteq \mathbb{C} $. Define $f(z)=z^2$. Define $g|U_1=z^2$ and $g|U_2=z^2+1$. Clearly $f$ is holoomorphic. Since $g|U_1,g|U_2$ are holomorphic, since $U_1 \cup U_2 \cong U_1 \amalg U_2$, and since coproducts of holomorphic functions are holomorphic, it follows that $g$ is holomorphic. Moreover, $\partial _zf=\partial _zg=2z$. However, they do not differ by a constant, but rather by the function
	\[
		z\mapsto 
		\begin{cases}
			0,&z\in U_1\\
			1,&z\in U_2
		\end{cases}
	\]
	which is only locally constant. Since $U$ is disconnected, this does not imply it is constant. Since disjoint connected components exist in $\mathbb{C} $ (consider two disjoint open balls), this is an example.
\end{proof}

\begin{problem}
	Let $f(z)=z^2$. Calculate that the integral of $f$ around the circle $\partial B_1(2)$ given by
	\[
		\int_{0}^{2\pi } f(2+e^{i\theta } ) \,\mathrm{d} \theta 
	\]
	is not zero. Yet the Cauchy theorem asserts that
	\[
		\oint_{\partial B_1(2)} f =0
	.\]
	Give an explanation.
\end{problem}
\begin{proof}
	Although $2+e^{i\theta } $ parameterizes $\partial B_1(2)$ counterclockwise for $\theta \in[0,2\pi ]$, the line integral of $f$ along $\gamma $ is defined as
	\[
		\int_\gamma f\mathrm{d} z = \int_0^{2\pi } \gamma ^*(f \mathrm{d} z)=\int_{0}^{2\pi } f(\gamma (\theta ))\gamma '(\theta ) \,\mathrm{d} \theta 
	.\]
	Note that the given integral is missing the factor of $\gamma '$. We compute the given integral:
	\[
		\int_0^{2\pi } 4+4e^{i\theta } +e^{2i\theta }  \,\mathrm{d} \theta =8\pi 
	\]
	since $e^{0} =e^{i2\pi } =e^{i4\pi } $.
\end{proof}

\begin{problem}
	\begin{enumerate}
		\item Let $\gamma $ be a positively oriented parameterization of $S^1$ with a counterclockwise orientation. Give an example of a $C^1$ function $f$ on a neighborhood of $\overline{B_1(0)} $ such that
			\[
				\int_\gamma f =0
			\]
			but such that $f$ is not holomorphic on any open set.
		\item Suppose that $f$ is $C^0$ on $B_1(0)$ and satisfies
			\[
				\int_{\partial B_r(0)} f=0
			\]
			for all $0<r<1$. Must $f$ be holomorphic on $B_1(0)$?
	\end{enumerate}
\end{problem}
\begin{proof}
	(1) Let $f(z)=\overline{z} ^2$, and let $z(\theta )=e^{i\theta } $ such that $\mathrm{d} z=ie^{i\theta } \mathrm{d} \theta $. Then
	\[
		\int_\gamma f = \int_0^{2\pi } e^{-2i\theta } ie^{i\theta }  \,\mathrm{d} \theta =-\int_{0}^{2\pi } e^{-i\theta }  \,\mathrm{d} \theta =0
	.\]
	However, clearly $\partial _{\overline{z} } f =2 \overline{z} \neq 0$ on all open $U\subseteq \mathbb{C} $.

	(2) Take $f$ the same as in (1), and parameterize $\partial B_r(0)$ by $z(\theta )=re^{i\theta } $. The integral is then
	\[
		\int_{0}^{2\pi } rie^{-i\theta }  \,\mathrm{d} \theta ,
	\]
	which is still 0.
\end{proof}



\begin{problem}
	Give a simple proof (actually a one-line proof if you use what we know about power series) that an analytic function has a \emph{local} primitive.
\end{problem}
\begin{proof}
	Any power series clearly has a primitive, and analytic functions admit local power series representations.
\end{proof}

\end{document}
